随着神经网络规模增大，单个 NPU 核心的计算能力、片上存储容量和核外 DDR 带宽成为相互制约的资源。本文研究给定有向无环计算图在同构多核 NPU 上的切图与调度问题：在满足数据依赖、核内容量、流水线和共享带宽约束的条件下，同时决定计算操作的子图归属、子图的核心分配以及各核心上的执行顺序。

本文以“驻留域”为统一建模对象。首先按子图或核心统计每份张量需要进入的驻留域，建立结构搬运量和 Task 激活约束；随后用四类有明确含义的起点生成候选，通过轻量时间骨架、官方核内展开和官方多核事件模拟逐层筛选。问题一针对每个子图独立成 Task 的场景，问题二允许同核子图合并并复用 L1/UB，问题三在问题二基础上加入共享只读 L2 Cache，并把消费者分布与查询先后纳入候选生成。

本稿当前为实验进行中的正文初稿。模型、公式、算法流程和验证协议已经写入；全量主矩阵、单核基线、消融、Cache 容量/带宽敏感性及其图表仍在运行，尚未写入任何实验数字。摘要中的关键结果位置统一保留为“\pending”，待对应运行目录完成、结果字段核对和队伍复核后再填入。本文不把程序运行状态、格式编译状态或历史试验结果当作本稿的科学结论。

\noindent\textbf{关键词：}多核 NPU；计算图切分；驻留域；片上存储；共享 Cache；离散事件模拟
